\section{Compositional Generative Modeling}
\label{sec:compositional_generative_modeling}

Modern generative models, particularly diffusion models, excel at capturing complex image 
distributions yet struggle when asked to produce \emph{novel semantic combinations} that were not explicitly 
observed during training \citep{liu2022compositional}.  
Even state-of-the-art diffusion models typically learn entangled representations in which global 
and local semantic factors---shape, texture, attributes, pathology type, spatial arrangement---interact nonlinearly.  
As a result, they cannot reliably synthesize images reflecting \emph{new concepts} assembled from previously 
learned factors, such as ``smiling woman with black hair’’ \citep{isometric_diffusion_icml_2024}
or clinically meaningful ``Silico--TB’’ combinations.  
This limitation reflects a broader challenge in generative modeling:  
\emph{current models lack explicit mechanisms for semantic factorization and composition}.

Compositional generative modeling seeks to overcome this limitation by enabling 
models to produce images that correspond to new combinations of latent factors in a structured way.  
The central goal is to construct generative mechanisms that respect 
the combinatorial structure of visual concepts.  
Earlier work in visually grounded imagination  
\citep{vedantam2017generative} formalizes this challenge through the “3C’s’’ criterion:  
\emph{correctness}, \emph{coverage}, and \emph{compositionality}.  
Similarly, hierarchical concept models such as SCAN  
\citep{Higgins2017SCANLH} demonstrate the need for structured latent spaces in which abstract 
and concrete visual concepts can be combined through controlled factor manipulations.

Despite impressive empirical demonstrations, the field still lacks a principled
theoretical account of when such combinations are semantically valid, why
certain operators fail, and which structural conditions enable meaningful hybrid
concepts.  Recent work \citep{mechanisms2025projective, du2024compositional} emphasise that compositional success is fundamentally governed by
latent alignment, disentanglement of semantic factors, and compatibility of
geometric structure across component models .

A coherent taxonomy emerging across the literature identifies \textbf{three
families of compositional operators}:  
(1) \textbf{Product-of-Experts (PoE) Composition},  
(2) \textbf{Mixture-Based Composition}, and     
(3) \textbf{Density-Based Composition (SuperDiff)}.  
These families differ in mathematical mechanism, assumptions, and the kinds of
semantic operations they can reliably express.  Diffusion models can be viewed
as a special case of energy-based models (EBMs) under the score-based
interpretation \citep{NEURIPS2020_49856ed4}, allowing PoE and mixture operators
to be applied directly to pretrained diffusion models.  Structural analyses
from \citet{mechanisms2025projective} provide a unifying theoretical
lens that clarifies when these operators succeed or break down.
\vspace{1em}
\subsubsection*{1. Product-of-Experts (PoE) Composition}

PoE is the classical approach to compositional generative modeling in both EBMs
and diffusion models. It constructs the intersection of multiple generative
components by multiplying their densities:
\[
p_{\mathrm{PoE}}(x) \propto \prod_{i=1}^k p_i(x),
\qquad
\log p_{\mathrm{PoE}}(x) = \sum_{i=1}^k \log p_i(x).
\]
Each expert imposes a constraint that the composite distribution must satisfy,
so PoE behaves as a logical AND operator. Using the score-based interpretation
of diffusion models \citep{NEURIPS2020_49856ed4}, this conjunction corresponds to
summing log-density gradients,
\[
\nabla_x \log p_{\mathrm{PoE}}(x)
    = \sum_{i=1}^k \nabla_x \log p_i(x),
\]
enabling PoE to be implemented directly within diffusion samplers without
retraining.

The practical difficulty lies not in the PoE operator but in the sampling
procedure. \citet{reduce_reuse_recycle_icml2023} show that naïve samplers fail even when the PoE
objective is correct, and they introduce MCMC-corrected Langevin samplers that
substantially improve stability and compositional quality. For PoE to succeed,
the experts must modulate distinct or conditionally independent generative
factors, share non-conflicting support, and exhibit compatible energy geometry.
The theoretical analysis of \citet{mechanisms2025projective} shows
that semantic correctness requires projective composition: latent projections of
each expert must interact linearly and preserve local independence within shared
tangent spaces. When these structural requirements are violated, PoE becomes
over-constrained, may collapse modes, produces veto effects where any
zero-density region forbids hybrid generation, and exhibits inconsistent
gradient directions leading to unstable visual artifacts. Despite its elegance,
PoE is therefore semantically brittle unless the underlying models possess
aligned and disentangled representations.

\vspace{1em}
% ==========================================================
\subsubsection*{2. Mixture-Based Composition}

Mixture-based composition forms the union of distributions rather than their
intersection. The composite density is given by
\[
p_{\mathrm{mix}}(x) = \frac{1}{Z} \sum_{i=1}^k \alpha_i p_i(x),
\]
corresponding to a logical OR operator. Unlike PoE, mixture operators expand the
set of plausible samples by accumulating density wherever \emph{any} expert
assigns high probability. As highlighted in \textit{Compositional Generative
Modeling: A Single Model is Not All You Need}, mixtures represent a distinct
semantic operation that cannot be implemented through score addition, because a
mixture of densities is not equal to a mixture of their score fields. Practical
sampling therefore requires explicit densities or weighted resampling schemes.

Mixtures are useful for modeling category variability, multimodal concepts, or
collections of plausible configurations. However, effective mixture-based
composition requires overlapping support regions, compatible noise schedules, and
density functions that do not produce contradictory high-probability regions.
Mixtures cannot represent conjunction and cannot enforce constraints between
attributes. They also blur incompatible modes, are sensitive to non-overlapping
supports, and require access to explicit densities—not natively provided by
score-based diffusion models. Thus, mixtures complement PoE operators by
providing a principled union mechanism, but they do not provide a mechanism for
composing fine-grained semantic factors into hybrid structures.

\vspace{1em}
% ==========================================================
\subsubsection*{3. Density-Based Composition (SuperDiff)}

Density-based composition, introduced by the SuperDiff framework
\citep{skreta2024superposition}, departs fundamentally from PoE and mixtures by
operating directly on the \emph{time-evolving densities} of pretrained diffusion
models. Given a shared forward SDE,
\[
\mathrm{d}x_t = f_t(x_t)\,\mathrm{d}t + g_t\,\mathrm{d}w_t,
\]
SuperDiff estimates the transient log-density $\log q_t(x_t)$ along the forward
noising trajectory. The key theoretical foundation is Theorem~1, which states:
\[
\frac{\mathrm{d}}{\mathrm{d}t}\log q_t(x_t)
  = -\mathrm{div}\, f_t(x_t)
    -\tfrac{1}{2} g_t^2 \|\nabla_x \log q_t(x_t)\|^2.
\]
SuperDiff provides an unbiased It\^o estimator for this expression that avoids
explicit computation of $\mathrm{div}\,f_t$, making density tracking feasible for
high-dimensional diffusion models.

By estimating intermediate densities, SuperDiff enables rigorous AND/OR
operators:
\[
\log q_{\mathrm{AND}} = \sum_i \lambda_i \log q_i, \qquad
\log q_{\mathrm{OR}} = \mathrm{softmax}(\{\log q_i\}).
\]
Compositions occur in density space rather than in score or energy space, giving
SuperDiff a stronger mathematical foundation. However, semantic correctness
still depends heavily on structural compatibility. Successful composition
requires a shared forward diffusion process, aligned latent geometry, local
independence of semantic factors, and compatible curvature (Hessian structure) of
the log-densities. These conditions match the projective composition framework of
\citet{mechanisms2025projective}. When violated, equal-density or
density-blending constraints can still produce semantically invalid hybrids
because density matching alone does not ensure semantic decomposability.
Nonlinear interactions between factors, curvature mismatch, and latent
misalignment all lead to artifacts or incoherent pathologies. SuperDiff is thus
a mathematically principled operator at the density level, but its semantic
reliability is similarly governed by structural alignment between the underlying
models.

\vspace{1em}
% ==========================================================
\subsection*{Summary}

Across PoE, mixtures, and density-level composition, a consistent principle
emerges: \emph{semantic composition is possible only when the component models
share aligned, disentangled, and geometrically compatible representations}. PoE
implements intersection-like conjunction, mixtures implement unions, and
SuperDiff implements mathematically rigorous density arithmetic. Each operator
has distinct strengths and limitations, but all exhibit predictable failure modes
when structural preconditions—in the sense established by
\citet{mechanisms2025projective}—are violated. These insights
directly motivate the research hypotheses of this thesis concerning latent
alignment, disentanglement, conditional factorisation, and curvature
compatibility as necessary conditions for clinically meaningful compositional
synthesis in medical diffusion models.